% =====================================================================
% Section 05, Phase 1 cleaning pipeline
% =====================================================================

\section{Phase 1 \textbullet\ Cleaning Pipeline}
\label{sec:phase1-cleaning}

The cleaning notebook (\fn{Cleaning Data.ipynb}) is a 12-section pipeline
that produces all 21 cleaned and merged CSV files used by the analysis
notebook. We walk through it section by section.

\subsection{Imports and CRS handling}

\begin{lstlisting}[caption={Cleaning notebook imports and CRS reprojection.}]
import os
import geopandas as gpd
import numpy as np
import pandas as pd
import plotly.graph_objects as go
from sklearn.impute import KNNImputer

CRS_METERS = "EPSG:32636" # UTM Zone 36N, metres for Cairo

boarding = boarding_raw.to_crs(CRS_METERS)
population = population_raw.to_crs(CRS_METERS)
speeds = speeds_raw.to_crs(CRS_METERS)
pax_flow = pax_flow_raw.to_crs(CRS_METERS)
routes = routes_raw.to_crs(CRS_METERS)
terminals = terminals_raw.to_crs(CRS_METERS)
veh_flow = veh_flow_raw.to_crs(CRS_METERS)
\end{lstlisting}

\subsection{Renaming and type casting}

The raw GeoJSON ships with abbreviated PTV-style column names
(\fn{ogc\_fid}, \fn{stop\_name}, \fn{daily\_all\_ba}, \ldots). The
notebook renames every column into a clean vocabulary
(\fn{Stop\_Name}, \fn{Daily\_All\_BA}, \ldots) so downstream code reads
naturally. Numeric fields are cast with \fn{pd.to\_numeric(\ldots,
errors="coerce")} so any unparseable strings become \fn{NaN} (rather
than crashing the pipeline). Nullable integer IDs use the pandas
\fn{Int64} dtype.

\subsection{KNN imputation, three targeted blocks}

\begin{callout}
\textbf{\color{plotblue}Why KNN, not simple imputation.}\quad The three
fields that need imputation (Boarding numeric counts, Routes \fn{Fare},
Pax-Flow \fn{Passengers\_Per\_Hour}) are \emph{spatially or
attributively structured}. A boarding count is correlated with what
nearby stops see. A microbus fare is correlated with route length and
vehicle type. Filling with column means or zeros would silently
flatten that structure. KNN with $k=5$, distance-weighted,
\fn{nan\_euclidean} metric uses the structure that's already there.
The notebook re-uses the same imputer policy for all three fields and
audits the result.
\end{callout}

\subsubsection{Boarding (spatial)}

For each row with at least one missing numeric, the imputer asks
``what do the five nearest stops in $(x, y, \text{numeric})$-space
look like?'' and fills with their distance-weighted average:

\begin{lstlisting}[caption={KNN imputation on Boarding stops, XY appended as features.}]
boarding["_x"] = boarding.geometry.x
boarding["_y"] = boarding.geometry.y
knn_features = boarding_num_cols + ["_x", "_y"]

imputer_boarding = KNNImputer(
 n_neighbors=5,
 weights="distance",
 metric="nan_euclidean",
)
X_imputed = imputer_boarding.fit_transform(
 boarding[knn_features].values
)

# write back only the cells that were null (clip to non-negative)
for col in boarding_num_cols:
 null_rows = boarding_null_mask[col]
 if null_rows.any():
 boarding.loc[null_rows, col] = (
 imputed_df.loc[null_rows, col].clip(lower=0)
 )
boarding.drop(columns=["_x", "_y"], inplace=True)
\end{lstlisting}

\subsubsection{Routes \fn{Fare} (attribute)}

The \stat{45.9\%}-null fare field is filled using route-attribute
features rather than spatial ones, because fare is governed by trip
length and vehicle type, not by where the route line happens to lie:

\begin{lstlisting}[caption={KNN imputation on Routes Fare, five attribute features.}]
fare_features = [
 "Fare", "Route_Length_km", "Vehicle_Capacity",
 "Direction", "_vehicle_code" # label-encoded vehicle type
]
imputer_fare = KNNImputer(n_neighbors=5, weights="distance",
 metric="nan_euclidean")
X_fare_imputed = imputer_fare.fit_transform(routes[fare_features].values)
routes.loc[fare_null_mask, "Fare"] = (
 fare_imp_df.loc[fare_null_mask, "Fare"].clip(lower=0).round(1)
)
\end{lstlisting}

\subsubsection{Passenger flow (segment-spatial)}

For passenger flow, the imputer uses each segment's centroid plus its
length, on the assumption that geometrically similar segments carry
similar throughput:

\begin{lstlisting}[caption={KNN imputation on Passenger Flow, centroid + length.}]
pax_flow["_cx"] = pax_flow.geometry.centroid.x
pax_flow["_cy"] = pax_flow.geometry.centroid.y
pax_flow["_seg_length"] = pax_flow.geometry.length
pax_features = ["Passengers_Per_Hour", "_cx", "_cy", "_seg_length"]

imputer_pax = KNNImputer(n_neighbors=5, weights="distance",
 metric="nan_euclidean")
X_pax_imputed = imputer_pax.fit_transform(pax_flow[pax_features].values)
pax_flow.loc[pax_null_mask, "Passengers_Per_Hour"] = (
 pax_imp_df.loc[pax_null_mask, "Passengers_Per_Hour"].clip(lower=0)
)
\end{lstlisting}

\subsection{Post-imputation null audit}

After all three KNN passes the cleaning notebook re-runs the audit. The
result is unambiguous: \stat{every numeric column across all 7 datasets is
null-free} (CLEAN) and \stat{zero rows were dropped} during the entire
process. This is the core property that lets us defend the merges and
the analysis: nothing was hidden, nothing was discarded.

\subsection{Engineered features}

Two computed columns are added to support downstream questions:

\begin{itemize}
 \item Boarding gets \fn{Total\_Morning\_Boarding},
 \fn{Total\_Evening\_Boarding}, and \fn{Total\_Routes\_At\_Stop}
, sums across the time-of-day breakouts.
 \item Routes gets a single classification field
 \fn{Transport\_Class} $\in$ \{\strval{Formal}, \strval{Informal}\}
 derived from vehicle type and capacity:
\end{itemize}

\begin{lstlisting}[caption={Formal/Informal classification on Routes.}]
is_informal = (
 routes["Vehicle_Type"].str.lower()
.str.contains("micro|coop|7-seat|7 seat", na=False)
 | (routes["Vehicle_Capacity"] <= 17)
)
routes["Transport_Class"] = np.where(is_informal, "Informal", "Formal")
\end{lstlisting}

This single column is the basis for every later question that compares
the two systems.
